\documentclass[12pt,a4paper]{article}

\usepackage[a4paper,margin=1in]{geometry}
\usepackage{graphicx}
\usepackage{booktabs}
\usepackage{amsmath,amssymb}
\usepackage{float}
\usepackage{hyperref}
\usepackage{xcolor}
\usepackage{listings}
\usepackage{enumitem}
\usepackage{fancyhdr}
\usepackage{longtable}

\hypersetup{colorlinks=true,linkcolor=blue,urlcolor=blue}

\pagestyle{fancy}
\fancyhf{}
\lhead{ICS1512}
\rhead{Machine Learning Algorithms Laboratory}
\cfoot{\thepage}

\lstset{
language=Python,
basicstyle=\ttfamily\small,
numbers=left,
frame=single,
breaklines=true,
showstringspaces=false
}

\begin{document}

\begin{tabular}{|p{4cm}|p{11cm}|}
\hline
Experiment No. & 2 \\ \hline
Experiment Title & Email Spam/Ham Classification using Na\"ive Bayes and KNN
\footnote{\textbf{GitHub Repository: }\url{https://github.com/Barath-j593/ML-assignments}}\\ \hline
Student Name & NavinKumar S \\ \hline
Register Number & 3122247001039\\ \hline
Faculty & Dr. Poreddy Ajay Kumar Reddy \\ \hline
Submission Date & 02/08/26 \\ \hline
\end{tabular}

\section{Objective}
To build spam/ham classifiers for the Spambase dataset using Na\"ive Bayes (Gaussian,
Multinomial, Bernoulli) and K-Nearest Neighbors, compare the three Na\"ive Bayes variants,
study the effect of varying $k$, compare the KDTree and BallTree neighbour-search algorithms,
optimize KNN using GridSearchCV and RandomizedSearchCV, evaluate computational complexity and
execution time, and evaluate all models using K-Fold cross-validation.

\section{Problem Statement}
Given the Spambase dataset of word/character frequency and capital-letter-run features
extracted from 4,601 emails, the problem is binary classification: predict whether each email
is spam (1) or not-spam (0). The input is a 57-dimensional numeric feature vector per email;
the output is a binary class label. The objective is to identify which Na\"ive Bayes variant
and which KNN configuration classify spam most accurately and efficiently, and to understand
the computational trade-offs between the two model families.

\section{Dataset Description}
\begin{longtable}{|p{6cm}|p{8cm}|}
\hline
Dataset Name & Spambase \\ \hline
Dataset Source & Kaggle -- \url{https://www.kaggle.com/datasets/somesh24/spambase} (UCI ML Repository) \\ \hline
Number of Samples & 4,601 \\ \hline
Number of Features & 57 (48 word-frequency, 6 character-frequency, 3 capital-run-length features) \\ \hline
Number of Classes & 2 (2,788 not-spam / 1,813 spam) \\ \hline
Missing Values & 0 \\ \hline
Train-Test Split & 80:20 stratified (3,680 train / 921 test), \texttt{random\_state=42} \\ \hline
\end{longtable}

\section{Exploratory Data Analysis}
Include all relevant visualizations.
\begin{itemize}
\item Dataset overview
\item Statistical summary
\item Missing value analysis
\item Class distribution
\item Correlation matrix
\item Feature distribution
\end{itemize}

\subsection*{Class Distribution}
\begin{figure}[H]
\centering
\includegraphics[width=0.5\linewidth]{class_distribution.png}
\caption{Class distribution -- Spambase}
\end{figure}
\textbf{Inference:} The dataset is moderately imbalanced (2,788 ham vs. 1,813 spam, about
61:39). This is mild enough that accuracy remains a reasonably informative metric, but
precision/recall/F1/ROC-AUC are still reported throughout to avoid overstating performance.
No missing values are present, so no missing-value visualization is required.

\subsection*{Correlation Matrix}
\begin{figure}[H]
\centering
\includegraphics[width=0.85\linewidth]{correlation_heatmap.png}
\caption{Correlation heatmap -- top 15 features most correlated with the spam label}
\end{figure}
\textbf{Inference:} Character frequency features such as \texttt{char\_freq\_!} and
\texttt{char\_freq\_\$}, along with word frequencies for ``your'', ``000'', ``remove'', and
``free'', show the strongest positive correlation with the spam label -- consistent with
common spam-email vocabulary and formatting (exclamation marks, dollar signs, promotional
words).

\subsection*{Feature Distribution}
\begin{figure}[H]
\centering
\includegraphics[width=0.48\linewidth]{hist_word_freq_free.png}
\includegraphics[width=0.48\linewidth]{hist_char_freq_excl.png}
\caption{Histograms of important features -- word\_freq\_free and char\_freq\_!}
\end{figure}
\textbf{Inference:} Both features are heavily right-skewed with a large mass at or near zero,
since most legitimate emails never use these tokens, while a minority of (mostly spam)
emails use them repeatedly -- exactly the long-tail pattern that makes them discriminative.

\begin{figure}[H]
\centering
\includegraphics[width=0.48\linewidth]{box_word_freq_free.png}
\includegraphics[width=0.48\linewidth]{box_char_freq_excl.png}
\caption{Boxplots of word\_freq\_free and char\_freq\_! grouped by class}
\end{figure}
\textbf{Inference:} Both features show a visibly higher median and wider spread for the spam
class (1) compared to ham (0), confirming they carry real discriminative signal rather than
noise -- this matches their high ranking in the correlation heatmap above.

\section{Data Preprocessing}
Explain:
\begin{itemize}
\item \textbf{Missing value handling:} Not required -- the dataset contains zero missing
values across all 4,601 rows and 57 features.
\item \textbf{Encoding:} Not required -- all 57 features and the target label are already
numeric.
\item \textbf{Feature scaling:} \texttt{StandardScaler} (zero mean, unit variance) was
applied for KNN and Gaussian NB, since distance-based and Gaussian-likelihood models are
sensitive to feature scale. Multinomial NB used the raw non-negative frequency values
directly (required, since it models non-negative counts). Bernoulli NB internally binarizes
each feature at a 0.0 threshold (scikit-learn default).
\item \textbf{Train-test split:} 80:20 stratified split (3,680 train / 921 test rows),
preserving the 61:39 class ratio in both subsets.
\item \textbf{Feature engineering:} Not applicable -- the 57 provided frequency/length
features were used as-is.
\end{itemize}

\section{Machine Learning Algorithm}
Explain:
\begin{itemize}
\item \textbf{Working principle:} Na\"ive Bayes applies Bayes' theorem,
$P(C\mid X) = P(X\mid C)P(C)/P(X)$, with a conditional independence assumption between
features, to compute each class's posterior probability directly from training-data
frequency/likelihood statistics -- there is no iterative optimization. Gaussian NB assumes
each feature is normally distributed within a class; Multinomial NB models features as
counts/frequencies; Bernoulli NB treats each feature as a binary presence/absence indicator.
KNN is a lazy, instance-based learner: it stores the training data and, at prediction time,
finds the $k$ nearest training points (by Euclidean/Manhattan distance) to a query point and
assigns the majority class among them (optionally weighted by inverse distance).
\item \textbf{Mathematical formulation:}
$$P(C\mid X) = \frac{P(X\mid C)\,P(C)}{P(X)}, \qquad
\hat{y} = \operatorname*{mode}_{i \in N_k(x)} y_i$$
where $N_k(x)$ is the set of $k$ nearest neighbours of query point $x$.
\item \textbf{Advantages:} Na\"ive Bayes is extremely fast to train and predict
($O(nd)$ / $O(d)$), works well even with limited data, and requires no hyperparameter tuning.
KNN makes no assumption about the underlying data distribution and can capture non-linear
decision boundaries.
\item \textbf{Limitations:} Na\"ive Bayes' independence assumption is frequently violated in
practice (many Spambase features are correlated, as seen in Section 4), which can bias its
probability estimates. KNN's prediction cost scales with training set size ($O(nd)$ for brute
force), it is sensitive to irrelevant/unscaled features, and it stores the entire training
set in memory.
\end{itemize}

\section{Experimental Setup}
\begin{tabular}{ll}
\toprule
Python Version & \\
Scikit-Learn Version & \\
IDE & Jupyter Notebook \\
Operating System & \\
Hardware & \\
\bottomrule
\end{tabular}

\section{Hyperparameter Selection}

\begin{tabular}{ll}
\toprule
Parameter & Value\\
\midrule
Best $n\_neighbors$ (GridSearchCV) & 9 \\
Best weights (GridSearchCV) & distance \\
Best metric (GridSearchCV) & manhattan \\
Best CV Accuracy (GridSearchCV) & 0.9253 \\
Best $n\_neighbors$ (RandomizedSearchCV) & 10 \\
Best weights (RandomizedSearchCV) & distance \\
Best metric (RandomizedSearchCV) & euclidean \\
Best CV Accuracy (RandomizedSearchCV) & 0.9234 \\
\bottomrule
\end{tabular}

\textbf{Note:} GridSearchCV exhaustively searched
$n\_neighbors \in \{1,3,5,7,9,11,13,15\} \times weights \in \{uniform, distance\} \times
metric \in \{euclidean, manhattan\}$ (32 combinations, 10.69s). RandomizedSearchCV sampled 20
combinations from $n\_neighbors \in [1,30]$ with the same weights/metric options (6.07s),
trading a marginal accuracy difference for roughly 1.75$\times$ less computation.

\section{Results}

\begin{tabular}{lc}
\toprule
Metric & Value (Best KNN, GridSearchCV-tuned)\\
\midrule
Accuracy & 0.9207 \\
Precision & 0.9421 \\
Recall & 0.8512 \\
F1-score & 0.8944 \\
ROC-AUC & 0.9712 \\
\bottomrule
\end{tabular}

\textbf{Comparison context:} The best Na\"ive Bayes variant, Bernoulli NB, achieved Accuracy
0.8762, Precision 0.8716, Recall 0.8044, F1 0.8367, ROC-AUC 0.9496 on the same test set --
Best KNN outperforms it on every metric except training/prediction speed (Section 9).

\section{Result Visualization}

Include all graphs generated during the experiment.

\begin{figure}[H]
\centering
\includegraphics[width=0.55\linewidth]{accuracy_vs_k.png}
\caption{Accuracy vs. k (KNN, unweighted Euclidean)}
\end{figure}
\textbf{Inference:} Accuracy rises steadily from $k=1$ to $k=11$ with no sign of turning
downward yet, suggesting the optimum $k$ for this dataset lies at or beyond $k=11$ --
consistent with GridSearchCV independently selecting $k=9$ once distance weighting and the
Manhattan metric are also allowed to vary.

\begin{figure}[H]
\centering
\includegraphics[width=0.75\linewidth]{gridsearch_heatmap.png}
\caption{GridSearchCV mean CV accuracy heatmap (n\_neighbors $\times$ weights, metric=euclidean)}
\end{figure}
\textbf{Inference:} Accuracy consistently improves moving from \texttt{uniform} to
\texttt{distance} weighting at every $k$, and generally increases with $k$ up to the tested
range -- both trends corroborate the accuracy-vs-$k$ curve above and the overall best-parameter
choice.

\begin{figure}[H]
\centering
\includegraphics[width=0.55\linewidth]{randomsearch_distribution.png}
\caption{RandomizedSearchCV mean CV score distribution across 20 sampled configurations}
\end{figure}
\textbf{Inference:} Most of the 20 randomly sampled hyperparameter combinations cluster in a
fairly narrow, high-accuracy band (roughly 0.90--0.925), showing KNN's performance on this
dataset is fairly robust to the exact choice of $k$/metric/weighting once a reasonable
neighbourhood size is used.

\begin{figure}[H]
\centering
\includegraphics[width=0.75\linewidth]{confusion_matrices.png}
\caption{Confusion matrices -- Gaussian NB, Multinomial NB, Bernoulli NB, Best KNN}
\end{figure}
\textbf{Inference:} Gaussian NB's confusion matrix visibly shows more false positives (ham
misclassified as spam) than any other model, matching its lower precision figure; Best KNN
shows the fewest total misclassifications of the four models.

\begin{figure}[H]
\centering
\includegraphics[width=0.48\linewidth]{roc_curves.png}
\includegraphics[width=0.48\linewidth]{pr_curves.png}
\caption{ROC curves and Precision-Recall curves for all four models}
\end{figure}
\textbf{Inference:} Best KNN's ROC curve dominates the others across almost the entire range
(highest AUC $=0.971$), with Bernoulli NB close behind ($0.950$). Multinomial NB's curve is
visibly the weakest, consistent with its lower accuracy and F1. The precision-recall curves
tell the same story more sharply, since PR curves are more sensitive to class imbalance than
ROC curves.

\begin{figure}[H]
\centering
\includegraphics[width=0.48\linewidth]{training_time_comparison.png}
\includegraphics[width=0.48\linewidth]{prediction_time_comparison.png}
\caption{Training time and prediction time comparison across all four models}
\end{figure}
\textbf{Inference:} All three Naive Bayes variants train and predict in single-digit
milliseconds, matching their theoretical $O(nd)$ training / $O(d)$ prediction cost. Best KNN
trains almost as fast as Naive Bayes but is roughly 70--150$\times$ slower at
\emph{prediction}, exactly matching the theoretical expectation that KNN prediction cost
scales with $n$ while NB prediction cost does not.

\begin{figure}[H]
\centering
\includegraphics[width=0.55\linewidth]{cv_accuracy.png}
\caption{5-fold cross-validation accuracy -- Bernoulli NB vs. best KNN}
\end{figure}
\textbf{Inference:} Best KNN outperforms Bernoulli NB on every single fold, with a
consistent margin of roughly 3--5 percentage points and noticeably lower fold-to-fold
variance, confirming its test-set advantage is stable rather than a lucky split.

\begin{figure}[H]
\centering
\includegraphics[width=0.7\linewidth]{classifier_comparison.png}
\caption{Classifier comparison -- Accuracy / Precision / Recall / F1}
\end{figure}
\textbf{Inference:} Best KNN leads on Accuracy, Precision, and F1, while Gaussian NB leads
only on Recall (at a real precision cost). Bernoulli NB is a clear second place overall,
ahead of both Gaussian and Multinomial NB on every metric except recall.

\section{Discussion}
KNN, once tuned, outperforms all three Na\"ive Bayes variants on every classification metric
for this dataset, at the cost of substantially higher prediction-time complexity -- a direct,
experimentally confirmed trade-off between accuracy and inference speed. Among the Na\"ive
Bayes variants, Bernoulli NB is the clear best choice, since Spambase's word-frequency
features carry most of their discriminative signal in simple presence/absence patterns
rather than exact counts or continuous distributions. KDTree and BallTree produced identical
accuracy (both are exact algorithms), with BallTree marginally faster at this scale;
GridSearchCV found a slightly better configuration than RandomizedSearchCV at roughly
1.75$\times$ the computation time. A limitation of this experiment is that the accuracy-vs-$k$
curve had not yet turned downward within the tested range, so the true single-parameter
optimum $k$ (holding weights/metric fixed) may lie beyond $k=11$; a wider sweep would clarify
this. Overfitting for very small $k$ (e.g. $k=1$, sensitive to individual noisy points) and
underfitting for very large $k$ (predictions collapsing toward the overall class balance) are
both theoretically expected, though this dataset's tested range did not clearly exhibit the
underfitting regime.

\section{Conclusion}
For the Spambase dataset, a properly tuned KNN classifier ($k=9$, distance-weighted, Manhattan
distance) is the best-performing model overall (Accuracy 0.9207, F1 0.8944, ROC-AUC 0.9712),
outperforming the best Na\"ive Bayes variant (Bernoulli NB) on every classification metric and
on every fold of 5-fold cross-validation. This advantage comes at a real computational cost:
KNN's prediction time is roughly two orders of magnitude slower than any Na\"ive Bayes
variant, which would become the dominant practical concern for a much larger deployment
dataset. Bernoulli NB is the preferred Na\"ive Bayes variant here, and GridSearchCV is
preferred over RandomizedSearchCV when the small additional computation cost is acceptable
for a marginally better hyperparameter configuration.

\section{References}
\begin{enumerate}[label={[\arabic*]}]
\item Scikit-learn Documentation, \url{https://scikit-learn.org/stable/}
\item A. G\'eron, \textit{Hands-On Machine Learning with Scikit-Learn, Keras and TensorFlow}, O'Reilly.
\item C. M. Bishop, \textit{Pattern Recognition and Machine Learning}, Springer.
\item Spambase Dataset, Kaggle, \url{https://www.kaggle.com/datasets/somesh24/spambase}
\end{enumerate}

\end{document}